\chapter{Probabiliteti dhe variablat e rastit}
\begin{qellime}
Studenti interpreton probabilitetin, shpërndarjen, pritjen dhe variancën; punon me ndryshore diskrete e të vazhdueshme; përdor ligjin e numrave të mëdhenj dhe teoremën qendrore kufitare; raporton gabimin statistik.
\end{qellime}

\section{Ngjarjet dhe probabiliteti}
Hapësira e mostrave $\Omega$ përmban rezultatet elementare. Probabiliteti plotëson $P(A)\ge0$, $P(\Omega)=1$ dhe aditivitetin për ngjarje të papajtueshme. Probabiliteti kushtor është
\begin{equation}
 P(A|B)=\frac{P(A\cap B)}{P(B)}.
\end{equation}
Pavarësia kërkon $P(A\cap B)=P(A)P(B)$. Ajo është supozim i fortë dhe nuk duhet ngatërruar me korrelacionin zero.

Teorema e Bayes-it,
\begin{equation}
 P(A|B)=\frac{P(B|A)P(A)}{P(B)},
\end{equation}
ndan informacionin paraprak nga informacioni i të dhënave. Në fizikën eksperimentale ajo shfaqet në klasifikim, zbulim sinjali dhe vlerësim parametrash.

\section{Ndryshoret e rastit}
Për ndryshore diskrete $X$, funksioni masor $p(x)$ plotëson $\sum_xp(x)=1$. Për ndryshore të vazhdueshme, dendësia $f(x)$ plotëson $\int f(x)\dd x=1$ dhe probabiliteti merret nga integrali. Funksioni kumulativ është $F(x)=P(X\le x)$.

Pritja dhe varianca janë
\begin{equation}
 \E[X]=\sum_xxp(x)\quad\text{ose}\quad\int xf(x)\dd x,
 \qquad \Var(X)=\E[(X-\E X)^2].
\end{equation}
Për një funksion $g$, $\E[g(X)]$ llogaritet mbi shpërndarjen e $X$; në përgjithësi $g(\E X)\ne\E[g(X)]$.

\section{Shpërndarje të rëndësishme}
Bernoulli përshkruan një provë me sukses $p$. Shuma e $n$ provave të pavarura është binomiale,
\begin{equation}
 P(K=k)=\binom nk p^k(1-p)^{n-k},\quad \E K=np,\quad \Var K=np(1-p).
\end{equation}

\begin{figure}[ht]
\centering\includegraphics[width=.74\textwidth]{14_shperndarja_binomiale.pdf}
\caption{Shpërndarja e numrit të stemave në 20 hedhje të një monedhe të paanshme.}
\end{figure}

Poisson përshkruan numërime të rralla me normë $\lambda$: $P(K=k)=e^{-\lambda}\lambda^k/k!$. Eksponencialja përshkruan kohën ndërmjet ngjarjeve Poisson dhe ka vetinë pa kujtesë. Gauss-i del nga mbledhja e shumë kontributeve të vogla; parametrat e tij janë mesatarja dhe varianca.

\section{Kovarianca dhe përhapja e pacaktueshmërisë}
\begin{equation}
 \Cov(X,Y)=\E[(X-\E X)(Y-\E Y)],\qquad
 \rho=\frac{\Cov(X,Y)}{\sigma_X\sigma_Y}.
\end{equation}
Për $z=f(\vect x)$ dhe pacaktueshmëri të vogla,
\begin{equation}
 \Var(z)\approx \nabla f^T\Sigma\nabla f,
\end{equation}
ku $\Sigma$ është matrica e kovariancës. Termat jashtë diagonales nuk mund të hiqen kur hyrjet janë të korreluara.

\section{Ligji i numrave të mëdhenj dhe CLT}
Mesatarja e mostrës
\begin{equation}
 \bar X_N=\frac1N\sum_{i=1}^NX_i
\end{equation}
afrohet te $\mu$. Për variancë të fundme, teorema qendrore kufitare jep
\begin{equation}
 \frac{\sqrt N(\bar X_N-\mu)}{\sigma}\Rightarrow\mathcal N(0,1).
\end{equation}
Gabimi standard është $s/\sqrt N$. Katërfishimi i mostrave përgjysmon gabimin, një ligj që përcakton koston e Monte Carlo-s.

\begin{figure}[ht]
\centering\includegraphics[width=.95\textwidth]{15_teorema_qendrore.pdf}
\caption{Shpërndarja e mesatares së ndryshoreve uniforme afrohet me formën Gauss kur rritet numri i termave.}
\end{figure}

\begin{kujdes}
Formula $s/\sqrt N$ kërkon mostra të pavarura ose një korrigjim për korrelacionin. Në seri Markov, madhësia efektive e mostrës mund të jetë shumë më e vogël se $N$.
\end{kujdes}

\begin{lidhje}
Laboratori simulon hedhjen e monedhës, ndërton intervale besimi dhe kontrollon shkallëzimin $N^{-1/2}$. Diskutohet dallimi ndërmjet luhatjes fizike dhe gabimit të vlerësuesit.
\end{lidhje}

\section*{Pyetje kontrolli}
\begin{enumerate}
\item Dalloni devijimin standard të popullatës nga gabimi standard i mesatares.
\item Kur kovarianca është thelbësore në përhapjen e pacaktueshmërisë?
\item Pse CLT nuk garanton konvergjencë të shpejtë për çdo shpërndarje?
\end{enumerate}
\section{Algoritmet bazë të simulimit probabilitar}
\subsection{Hedhja e monedhës}
\begin{lstlisting}[language=Python,caption={Vlerësimi i probabilitetit dhe gabimit standard.}]
import numpy as np

def monedha(N, p=0.5, seed=2026):
    rng = np.random.default_rng(seed)
    x = rng.random(N) < p
    p_hat = x.mean()
    se = np.sqrt(p_hat*(1-p_hat)/N)
    return p_hat, se, x

for N in [100, 1_000, 10_000, 100_000]:
    p_hat, se, _ = monedha(N)
    print(N, p_hat, se)
\end{lstlisting}
Vlerësimi duhet përsëritur për shumë fara për të parë shpërndarjen e $\hat p$, jo vetëm një trajektore të vetme.

\subsection{Intervali Wilson}
Intervali normal $\hat p\pm z\sqrt{\hat p(1-\hat p)/N}$ sillet dobët për $N$ të vogël ose $\hat p$ pranë 0 e 1. Intervali Wilson është më i qëndrueshëm:
\begin{equation}
 \frac{\hat p+z^2/(2N)\pm z\sqrt{\hat p(1-\hat p)/N+z^2/(4N^2)}}{1+z^2/N}.
\end{equation}

\begin{lstlisting}[language=Python]
def interval_wilson(k, n, z=1.96):
    phat = k/n
    den = 1 + z*z/n
    q = phat + z*z/(2*n)
    d = z*np.sqrt(phat*(1-phat)/n + z*z/(4*n*n))
    return (q-d)/den, (q+d)/den
\end{lstlisting}

\subsection{Përhapja Monte Carlo e pacaktueshmërisë}
\begin{lstlisting}[language=Python,caption={Përhapja e hyrjeve të korreluara.}]
rng = np.random.default_rng(7)
mu = np.array([10.0, 2.0])
Sigma = np.array([[0.25, 0.08], [0.08, 0.04]])
x = rng.multivariate_normal(mu, Sigma, size=200_000)
z = x[:, 0] / x[:, 1]**2
print(np.mean(z), np.std(z), np.quantile(z, [0.025, 0.975]))
\end{lstlisting}
Krahasimi me linearizimin $\nabla f^T\Sigma\nabla f$ tregon kur jolineariteti dhe asimetria bëhen të rëndësishme.

\subsection{Kontrolli i CLT-së}
Për çdo $N$ gjenerohen shumë eksperimente të pavarura dhe studiohet
\begin{equation}
 Z=\frac{\sqrt N(\bar X-\mu)}{\sigma}.
\end{equation}
Histogrami i $Z$ krahasohet me Gauss-in standard. Termi i rekomanduar është \emph{shpërndarje normale} ose \emph{Gauss}, jo përkthime jo të qëndrueshme.

\begin{kujdes}
\emph{Luhatje} përshkruan variacionin e një madhësie të rastit; \emph{gabim standard} përshkruan pacaktueshmërinë e një vlerësuesi. Ato nuk janë sinonime.
\end{kujdes}

Për probabilitetin, statistikën dhe terminologjinë e vlerësimit janë konsultuar materiale universitare shqip \cite{capollari,butka,tasho}; lidhja me simulimin fizik mbështetet te \cite{newman,landau}.

\subsection{Nga aksiomat te ndryshorja e rastit}
Një hapësirë probabilitare $(\Omega,\mathcal F,P)$ përbëhet nga rezultatet e mundshme, ngjarjet e matshme dhe masa e probabilitetit. Ndryshorja e rastit $X:\Omega\to\R$ e shndërron rezultatin elementar në madhësi numerike. Për një ndryshore diskrete,
\begin{equation}
 \E[g(X)]=\sum_x g(x)p_X(x),
\end{equation}
ndërsa për një ndryshore të vazhdueshme,
\begin{equation}
 \E[g(X)]=\int_{-\infty}^{\infty}g(x)f_X(x)\dd x.
\end{equation}
Varianca del nga identiteti
\begin{equation}
 \Var(X)=\E[(X-\mu)^2]=\E[X^2]-\mu^2.
\end{equation}
Ajo ka njësinë në katror të $X$; devijimi standard ka të njëjtën njësi si madhësia fizike.

\subsection{Hedhja e monedhës dhe shpërndarja binomiale}
Le të jetë $X_i\in\{0,1\}$ treguesi i suksesit me $P(X_i=1)=p$. Për prova të pavarura, $S_N=\sum_iX_i$ ka shpërndarje binomiale. Numri i vargjeve me $k$ suksese është $\binom Nk$, prandaj
\begin{equation}
 P(S_N=k)=\binom Nk p^k(1-p)^{N-k}.
\end{equation}
Nga lineariteti i pritjes dhe pavarësia,
\begin{equation}
 \E[S_N]=Np,\qquad \Var(S_N)=Np(1-p).
\end{equation}
Vlerësuesi $\widehat p=S_N/N$ është i paanshëm dhe ka variancë $p(1-p)/N$. Kështu gabimi standard ulet si $N^{-1/2}$, jo si $N^{-1}$.

\subsection{Ligji i numrave të mëdhenj dhe teorema qendrore kufitare}
Ligji i numrave të mëdhenj pohon se mesatarja e mostrës $\overline X_N$ afrohet me $\mu$ kur $N$ rritet. Teorema qendrore kufitare jep më shumë: për ndryshore të pavarura me variancë të fundme,
\begin{equation}
 \frac{\sqrt N(\overline X_N-\mu)}{\sigma}
 \xrightarrow{d}\mathcal N(0,1).
\end{equation}
Kjo justifikon intervalin asimptotik $\overline X\pm z_{1-\alpha/2}s/\sqrt N$. Për mostra të vogla me variancë të panjohur përdoret shpërndarja Student. Për Bernoulli me $p$ pranë zeros ose njësisë, intervali normal mund të dalë jashtë $[0,1]$; intervali Wilson ka sjellje më të mirë.

\subsection{Kovarianca dhe përhapja e pacaktueshmërisë}
Për vektorin $\vect X$ me kovariancë $\Sigma$ dhe $Y=g(\vect X)$, linearizimi rreth mesatares jep
\begin{equation}
 g(\vect X)\simeq g(\vect\mu)+\nabla g(\vect\mu)^T(\vect X-\vect\mu),
\end{equation}
prandaj
\begin{equation}
 \Var(Y)\simeq\nabla g(\vect\mu)^T\Sigma\nabla g(\vect\mu).
\end{equation}
Termat jashtë diagonales mund ta rrisin ose ta ulin pacaktueshmërinë. Supozimi i pavarësisë pa argument mund të jetë po aq problematik sa neglizhimi i një burimi gabimi.

Kur linearizimi nuk është i mjaftueshëm, përhapja Monte Carlo gjeneron $\vect X^{(r)}$ nga shpërndarja e hyrjeve dhe llogarit $Y^{(r)}=g(\vect X^{(r)})$. Shpërndarja empirike e $Y$ mund të jetë asimetrike; atëherë raportohen kuantilet dhe jo vetëm mesatarja me devijimin standard.

\begin{lstlisting}[language=Python,caption={Përhapje Monte Carlo e hyrjeve të korreluara.}]
import numpy as np

rng = np.random.default_rng(2026)
mes = np.array([2.0, 3.0])
kov = np.array([[0.04, 0.03],
                [0.03, 0.09]])
x = rng.multivariate_normal(mes, kov, size=100_000)
y = x[:, 0]**2 / x[:, 1]
q025, q50, q975 = np.quantile(y, [0.025, 0.5, 0.975])
print(q50, q025, q975)
\end{lstlisting}

\subsection{Përmbledhje dhe ushtrime}
Probabiliteti përshkruan modelin e rastësisë; statistika përdor të dhënat për të nxjerrë përfundime rreth këtij modeli. Pacaktueshmëria e rezultatit duhet lidhur me një procedurë kampionimi dhe me supozimet e saj.
\begin{enumerate}
\item Nxirrni mesataren dhe variancën e binomiales nga funksioni gjenerues i momenteve.
\item Krahasoni intervalin normal dhe Wilson për $k=1$ sukses në $N=20$ prova.
\item Provoni me simulim shkallëzimin $N^{-1/2}$ dhe shpjegoni shpërndarjen e pjerrësive të matura.
\item Për periodën $T=2\pi\sqrt{L/g}$, nxirrni përhapjen lineare të pacaktueshmërisë dhe krahasojeni me Monte Carlo.
\end{enumerate}

\subsection{Vlerësuesi, animi dhe gabimi kuadratik mesatar}
Një vlerësues $\widehat\theta$ është ndryshore rasti sepse varet nga mostra. Animi është
\begin{equation}
 \operatorname{Bias}(\widehat\theta)=\E[\widehat\theta]-\theta.
\end{equation}
Gabimi kuadratik mesatar ndahet si
\begin{equation}
 \E[(\widehat\theta-\theta)^2]
 =\Var(\widehat\theta)+\operatorname{Bias}(\widehat\theta)^2.
\end{equation}
Ky identitet tregon se vlerësuesi i paanshëm nuk është automatikisht optimal: një anim i vogël mund të pranohet nëse ul ndjeshëm variancën. Në simulimet fizike, kjo shfaqet te rregullimi i modeleve, prerja e peshave ekstreme ose vlerësuesit me kontroll variabël.

Varianca e mostrës me emërues $N-1$ është e paanshme për variancën e popullatës kur mostrat janë të pavarura,
\begin{equation}
 s^2=\frac{1}{N-1}\sum_{i=1}^{N}(X_i-\overline X)^2.
\end{equation}
Emëruesi $N-1$ vjen sepse mbetjet kanë shumën zero dhe vetëm $N-1$ shkallë lirie. Për seri të korreluara, formula ende përshkruan shpërndarjen margjinale, por $s/\sqrt N$ nuk është gabimi standard i mesatares.

\subsubsection{Intervalet e besimit dhe interpretimi}
Një interval besimi $95\%$ nuk do të thotë se, pasi të dhënat janë vëzhguar, parametri fiks ka probabilitet $95\%$ të jetë brenda atij intervali. Interpretimi frekuentist është se procedura prodhon intervale që mbulojnë parametrin në $95\%$ të përsëritjeve ideale. Ky dallim është i rëndësishëm në raportim.

Për një mesatare me variancë të panjohur dhe të dhëna Gauss të pavarura,
\begin{equation}
 \frac{\overline X-\mu}{s/\sqrt N}\sim t_{N-1}.
\end{equation}
Për mostra të mëdha, kuantilet $t$ afrohen me ato normale. Për probabilitete binomiale, intervali Wilson ose një interval i bazuar në gjasë është më i qëndrueshëm se intervali Wald, sidomos pranë kufijve.

\subsubsection{Bootstrap-i}
Bootstrap-i joparametrik trajton shpërndarjen empirike si përafrim të popullatës. Nga mostra origjinale me $N$ elemente krijohen mostra të reja me rizgjedhje dhe të njëjtën madhësi. Për secilën llogaritet statistiku $T^{*(b)}$. Shpërndarja e këtyre statistikëve vlerëson variancën, animin dhe intervalet.

Metoda funksionon mirë për shumë statistikë të lëmuar, por nuk është universale. Për maksimumin, bishtat ekstremë dhe mostra shumë të vogla mund të jetë e dobët. Për seri kohore nuk duhet rizgjedhur çdo pikë veçmas, sepse prishet varësia; përdoret bootstrap në blloqe. Në Monte Carlo me peshë, skema e rizgjedhjes duhet të respektojë mënyrën si u krijuan peshat.

\subsubsection{Kovarianca nga burime të përbashkëta}
Në eksperimente fizike, dy madhësi shpesh varen nga i njëjti kalibrim. Nëse $x=x_0(1+\epsilon)$ dhe $y=y_0(1+\epsilon)$ ndajnë të njëjtin gabim relativ, raporti $x/y$ anulon gabimin e përbashkët në rendin e parë. Trajtimi i $x$ dhe $y$ si të pavarura do ta mbivlerësonte pacaktueshmërinë. Në diferencën $x-y$, i njëjti korrelacion mund ta ulë ose ta rrisë efektin në varësi të shenjave.

Matrica e kovariancës duhet të jetë simetrike pozitivisht gjysmë e përcaktuar. Një matricë e ndërtuar nga koeficientë korrelacioni të papajtueshëm mund të ketë vlera vetjake negative dhe të mos përfaqësojë asnjë shpërndarje. Para kampionimit shumënormal kontrollohen simetria dhe spektri; vlera të vogla negative nga rrumbullakimi mund të kërkojnë korrigjim të dokumentuar.

\subsubsection{Ligji total i variancës}
Kur modeli varet nga parametri i rastit $\Theta$ dhe nga zhurma e brendshme,
\begin{equation}
 \Var(Y)=\E[\Var(Y\mid\Theta)]
 +\Var(\E[Y\mid\Theta]).
\end{equation}
Termi i parë mat variabilitetin brenda realizimeve për parametra të fiksuar; i dyti mat ndryshimin ndërmjet parametrave. Kjo ndarje është e dobishme në simulimin e zbritjes së mjetit hapësinor, ku mund të dallohen ndryshimet atmosferike nga zhurma e sensorit ose e kontrollit.

